% =========================
% Cas pratique 12
% =========================

\section[Cas 12]{Cas pratique 12 : Standard Essential Patent et
abus de position dominante}

%--------------------------
%--------------------------
\begin{frame}{Question}

\begin{block}{}
\begin{itemize}
  \item \textbf{Art. 102 TFUE} (abus de position dominante).
  \item \textbf{CJUE, 16 juillet 2015, Huawei Technologies / ZTE} (SEP + FRAND).
\end{itemize}
\end{block}

Se comporter comme un \textbf{willing licensee} (candidat sérieux à une licence FRAND),
afin de rendre une demande d'injonction \textbf{risquée / abusive} pour le titulaire du SEP
(standard essential patent).

\end{frame}

%--------------------------
\begin{frame}{Question}

Trois conditions d'un "wiling licensee":
\begin{itemize}
  \item \textbf{Réagir sans délai} dès la notification/ mise en demeure : manifester
  une volonté claire d'obtenir une licence \textbf{FRAND} (pas de stratégie dilatoire).
  
  \item \textbf{Négocier de bonne foi} : demander une offre FRAND, échanger sur la base
  (portefeuille, périmètre, taux, assiette, pays, durée), coopérer activement.
  
  \item \textbf{Si désaccord} : formuler une \textbf{contre-offre écrite} à des conditions FRAND
  et fournir une \textbf{garantie} (séquestre / garantie bancaire) pour les redevances,
  avec \textbf{reddition de comptes} des actes d'exploitation.
\end{itemize}

L'interdiction est évitable si TELEPHONICA
met en oeuvre les trois mesures préventives.

\end{frame}